\chapter[Constructing References, Leaving Footnotes, and Copyright]%
{%
  \label{ch:make-bib}%
  Constructing References, Leaving Footnotes, and Copyright\footnotemark
}
% the footnote text goes after the heading so that the ToC and \label point to the chapter
\footnotetext{%
  \noindent\hangindent=1.5em\hangafter=1%
  Chapter 5 was previously published and is reproduced with permission from \textcite{author_year}. [Note: This is an example of permission to reproduce material from a book or journal article, with the full reference corresponding to the narrative citation (Author, [Year]) given in the References.]%
}
Your References section is automatically generated using \texttt{biblatex}.
New \texttt{BibTeX} entries can be added to the file \texttt{ref.bib} and are only included in the final References if they are actually \texttt{cite}d in your document using the \texttt{\backslash cite} command.
For example:
``We follow \cite{evans_pde}\ldots''.
The specific formatting of the References depends on your discipline; for this reason, these options can be changed inside your \texttt{preamble.tex} file.
Here, we are using \textcite{lax_hyperbolic} as an example narrative citation.
As an example parenthetical citation, here is a reference on the interplay of parasite mating systems and complex life cycle evolution \parencite{hulke_criscione}.
The References section should be double spaced between references; single-spacing within a reference that runs over one line is optional.
Anywhere in the document, you may leave footnotes using the \texttt{\backslash footnote} command.\footnote{%
  \noindent\hangindent=1.5em\hangafter=1 This is an example of an in-text footnote if you need to indicate an aside in your main text or indicate permission for other copyrighted material.%
}
This command can be used for citations,\footnote{%
  \Autocite{lax_hyperbolic}.%
}
which will be formatted according to your \texttt{biblatex} settings.

\section{Bibliography Styles}
In \texttt{preamble.tex}, we set up the bibliography using:
\begin{verbatim}
      \usepackage[... style=authoryear ... ]{biblatex}
\end{verbatim}
However, this style is not mandatory and should be changed as appropriate for your discipline.
The \texttt{biblatex} package supports a wide range of standard styles, including numeric, alphabetic, APA, MLA, IEEE, Nature, \textit{etc}.
All of these options are detailed in \texttt{biblatex}'s
\href{https://ctan.org/pkg/biblatex?lang=en}{documentation on CTAN} and more concisely in this
\href{https://tug.ctan.org/info/biblatex-cheatsheet/biblatex-cheatsheet.pdf}{\texttt{biblatex} cheat-sheet}.
Of particular importance are the \texttt{style} and \texttt{citestyle} options, where the latter is used to style in-text citations in particular.
One may also use the \href{https://ctan.org/pkg/amsrefs}{\texttt{amsrefs}} package to format their references in conjunction with \texttt{biblatex} or as a replacement.
The \texttt{cleveref} package is only \href{https://latex3.github.io/tagging-project/tagging-status/}{partially compatible with tagging}  and should be used with caution at this time.

\section{Copyright}
Please refer to the TAMU \href{https://grad.tamu.edu/degree-completion/thesis-and-dissertation/writing-tools/index.html}{Guidelines for Theses, Dissertations, and Records of Study} for policies regarding the use of your previously published material or others’ copyrighted material in your manuscript.

Proper acknowledgment and permission statements must be included in the manuscript, with formatting depending on the type and extent of material used: (1) for previously published chapters or substantial portions of a chapter, include a footnote at the end of the chapter title, as shown in the Chapter~5 title; (2) for limited reproduced text, include a footnote at the end of the relevant sentence (as shown for footnote 2 on the prior page); and (3) for reproduced tables and figures, include a “reproduced with permission” statement as the final sentence of the table caption or figure legend, along with a narrative citation (Author, [Year]), and do not use a footnote.
